
En la actualidad, la carga computacional ha migrado del dispositivo local hacia grandes Centros de Procesamiento de Datos (CPD), los cuales emplean una arquitectura de microservicios basada
en contenedores y gestionada por orquestadores como Kubernetes. Esta cambio de paradigama, a permitido aumentar la capacidad de los servicios y aplicaciones para dar servicios a multiples nuevos
clientes, asi como la ejecucion simultanea de tareas de alta computacion en paralelo, las cuales antes estaban limitadas por la capacidad hardware del proveedor. Como consecuencia, el consumo de recursos
se ha visto altamnete afectado, siendo uno de los mas afectados el consumo energetico, el cual segun se indica en \cite{ec-cloud-worflowscheduling} supone entre el 15 y 20\% del consumo global energetico.\

Siendo mas precisos, basandome en el articulo \cite{survery-videoencoding} sabemos que la navegación por internet supone un 3.7\% de los Gases de Efecto Invernadero, del cual el 65\% corresponde al 
streaming de video. Esto supone que el stremaing de video es responsable de las emisiones de efecto invernadero equivalentes a las liberadas por la industra de transporte aheros, como se cita el paper
\cite{survery-videoencoding}.\

Es por este motivo que la obtención de métricas precisas, segmentadas por proceso y por area de hardwared utilizada, se ha convertido en una necesidad vital. Tanto en el ambito de la estimacion 
de usos y costes, como en la gestion de analisis de sostenibilidad, siendo estos ultimos esenciales para la toma de decisiones a futuro, tanto en el contexto industrial como el contexto social\/politico

% Contexto Del Trabajo
Bajo este contexto, se desarrolla el siguiente Trabajo de Fin de Máster (TFM), siguiendo una de las lineas de investigacion del proyecto Energy-Conscious Optimization for Cloud-Edge Video Ecosystems 
(ECO-STREAM). Este proyecto tiene como objetivo la optimización de cargas de trabajo de codificación y streaming de vídeo de alta calidad, en 
arquitectura de contenedores, mediante el uso de CPU y GPU físicas.

El objetivo principal de este trabajo de fin de mates es 
definir la configuración de codificación de alta resolución que minimice el consumo energético sin comprometer la calidad del resultado final. Para ello, 
es indispensable contar con herramientas capaces de medir el consumo energético a nivel de contenedor de forma precisa, segmentando el gasto entre CPU y GPU.\

Para la adquisición de métricas, se propone el uso de herramientas software adaptadas para el contexto de las cargas de trabajo, utilizando la medición de consumo energético total del
sistema como punto de comparacion y verificación. El uso de herramientas hardware, se ha descartado debido a su complejidad de implementación, así como
a los desafíos de mantenimiento y escalabilidad que presentan en entornos de producción.\

Las herramientas seleccionadas utilizarán Running Average Power Limit Energy Reporting (RAPL) \cite{rapl-tool},
una aplicacion de intel la cual permite extraer metricas precisas del consumo energetico de varios componentes dle sistema, junto con la api interna de NVIDIA y sus librerias asociadas.
Cabe destacar que RAPL agrupa en grupos y subgrupos los componentes del sistema, como se observa en la imagen \ref{fig:rapl-too-working}, 
permitiendo obtener metricas precisas. 

% Resumen e indicacion de los capitulos
Una vez definido el contexto y la problemática que se desea solventar en la introducción \ref{intro}, pasaremos al capítulo \ref{mot-obj} (Motivacion y Objetivos)
donde se remarcaran los objetivos y motivaciones del
proyecto. Posteriormente, en el capitulo \ref{estado-del-arte} (Estado del Arte) se definirá el marco teórico y contexto tecnológico, el cual será la base con la que justificaremos las decisiones tomadas a lo largo 
de este trabajo. Seguidamente, en el capítulo \ref{metodo} (Desarrollo del Proyecto) se describirán los diferentes pasos, 
configuraciones y problemáticas encontradas durante el desarrollo del trabajo. En el capitulo  \ref{results} (Resultados y Pruebas), se hara uso los datos y resultados obtenidos del apartado 
anterior y se agregarán diferentes tablas y ayudas gráficas las cuales facilitarán la comprensión y el desarrollo de las conclusiones. Para finalizar, en el capítulo \ref{conclusion} (Conclusiones y Trabajo Futuro), 
se expondrán las conclusiones y reflexiones acerca del trabajo realizado. Destacar que al final de este trabjo tambien se podra encontrar la bibliografia \ref{biblio}   y un anexo \ref{apendice} , el cual 
contendrá informacion adicional.    
